\documentclass[10pt]{beamer}
\usetheme{metropolis}
\usepackage{graphicx}  % 插入图片的核心宏包，必须加！
\usepackage{amsmath}
\usepackage{booktabs}
\usepackage{physics}
\usepackage{xcolor}
\usepackage{tikz}
\usetikzlibrary{quantikz}
\usepackage{wasysym} % 提供\yes（勾）和\no（叉）命令



% 设置颜色
\definecolor{SUSTechBlue}{RGB}{0, 82, 155}
\setbeamercolor{structure}{fg=SUSTechBlue}

% 标题信息
\title{E91 Quantum Cryptography Protocol}
\subtitle{A Simulation Study Based on Bell's Theorem}
\author{Zhao Yumu}
\institute{School of Physical Science and Technology \\ ShanghaiTech University}
\date{24/12/2025}
\titlegraphic{\includegraphics[width=2cm]{ShanghaiTech_logo.png}}

\begin{document}

% ============ 标题页 ============
\begin{frame}
    \titlepage
\end{frame}

% ============ 目录 ============
\begin{frame}{Presentation Outline}
    \tableofcontents
\end{frame}

% ============ 第一部分：背景介绍 ============
\section{Background \& Motivation}

\begin{frame}{The Challenge in Classical Cryptography}
    \begin{itemize}
        \item \textbf{Key Distribution Problem}: How to securely share secret keys?
        \item Classical channels can be passively monitored
        \item Mathematical security ≠ Physical security
        \item Need for \alert{information-theoretically secure} protocols
    \end{itemize}
    
    \vspace{5mm}
    \begin{block}{Quantum Solution}:Quantum mechanics provides fundamentally new security guarantees
    \end{block}
\end{frame}

\begin{frame}{Bell's Theorem \& Quantum Nonlocality}
    \begin{columns}
        \begin{column}{0.5\textwidth}
            \begin{itemize}
                \item Bell's inequality (1964)
                \item Tests local hidden variable theories
                \item Violation implies quantum nonlocality
                \item CHSH inequality: $|S| \leq 2$
            \end{itemize}
        \end{column}
        \begin{column}{0.5\textwidth}
            \begin{figure}
                \centering
                \includegraphics[width=\textwidth]{bell_test.png}
                \caption{Bell test experiment}
            \end{figure}
        \end{column}
    \end{columns}
    
    \[
    S = E(a_1, b_1) - E(a_1, b_3) + E(a_3, b_1) + E(a_3, b_3)
    \]
    Quantum prediction: $|S| = 2\sqrt{2} \approx 2.828$
\end{frame}

% ============ 第二部分：E91协议原理 ============
\section{E91 Protocol Principles}

\begin{frame}{E91 Protocol Overview (Ekert, 1991)}
    \begin{columns}
        \begin{column}{0.6\textwidth}
            \begin{itemize}
                \item Uses entangled photon pairs
                \item Security based on Bell's theorem
                \item No key information during transmission
                \item Key "comes into being" after measurement
            \end{itemize}
            
            \vspace{1mm}
            \textbf{Three measurement bases:}
            \begin{itemize}
                \item Alice: a_1=$0^\circ, a_2=45^\circ, a_3=90^\circ$
                \item Bob: b_1=$45^\circ, b_2=90^\circ, b_3=135^\circ$
            \end{itemize}
       \begin{align*}
E(a_i, b_j) &= P_{++}(a_i, b_j) + P_{--}(a_i, b_j) \\
&\quad - P_{+-}(a_i, b_j) - P_{-+}(a_i, b_j)
\end{align*}
       \[
    S = E(a_1, b_1) - E(a_1, b_3) + E(a_3, b_1) + E(a_3, b_3)
    \]

        \end{column}
        
        \begin{column}{0.4\textwidth}
            \begin{figure}
                \centering
                \includegraphics[width=\textwidth]{e91_setup.png}
                \caption{E91 experimental setup}
            \end{figure}
        \end{column}
    \end{columns}
\end{frame}

\begin{frame}{Protocol Workflow}
    \begin{enumerate}
        \item \textbf{Source} emits singlet state pairs: $\frac{|01\rangle - |10\rangle}{\sqrt{2}}$
        \item \textbf{Alice \& Bob} randomly choose measurement bases
        \item \textbf{Public discussion} of basis choices
        \item \textbf{Bell test} using mismatched bases
        \item \textbf{Key generation} from matched bases
        \item \textbf{Privacy amplification} \& error correction
    \end{enumerate}
    
    \vspace{5mm}
    \begin{alertblock}{Security Check}
        If $|S| > 2$, channel is secure. If $|S| \leq 2$, eavesdropping detected!
    \end{alertblock}
\end{frame}

\begin{frame}{Quantum State Preparation}
    \centering
    \begin{quantikz}[transparent]
        \lstick{$\ket{0}$} & \gate{H} & \ctrl{1} & \gate{Z} & \meter{} \\
        \lstick{$\ket{0}$} & \qw & \targ{} & \gate{X} & \meter{}
    \end{quantikz}
    
    \vspace{5mm}
    \begin{itemize}
        \item Initial state: $|00\rangle$
        \item Apply Hadamard + CNOT + phase flip
        \item Final singlet state: $\frac{|01\rangle - |10\rangle}{\sqrt{2}}$
        \item Perfect anticorrelation for same measurement basis
    \end{itemize}
    \begin{equation*}
        E(\mathbf{a}_i, \mathbf{b}_j) = -\mathbf{a}_i \cdot \mathbf{b}_j = -\cos(\theta_a - \theta_b)
    \end{equation*}
    
\end{frame}

% ============ 第三部分：模拟实验设计 ============
\section{Simulation Design}

\begin{frame}{Python Simulation Framework}
    \begin{columns}
        \begin{column}{0.5\textwidth}
            \textbf{Key Components:}
            \begin{itemize}
                \item \texttt{E91Simulator} class
                \item Qiskit for quantum circuits
                \item NumPy for numerical calculations
                \item Matplotlib for visualization
            \end{itemize}
            
            \vspace{3mm}
            \textbf{Two scenarios:}
            \begin{enumerate}
                \item Ideal (no eavesdropping)
                \item With Eve's intercept-resend attack
            \end{enumerate}
        \end{column}
        
        \begin{column}{0.5\textwidth}
            \begin{block}{Simulation Parameters}
                \begin{itemize}
                    \item $N = 10,000$ photon pairs
                    \item Measurement bases as in paper
                    \item Eve attacks 100\% of transmissions
                    \item Random seed for reproducibility
                \end{itemize}
            \end{block}
        \end{column}
    \end{columns}
\end{frame}

\begin{frame}{Eavesdropper Model}
    \textbf{Eve's intercept-resend attack:}
    \begin{enumerate}
        \item Intercepts photon going to Bob
        \item Measures in random basis
        \item Sends new photon to Bob in measured state
        \item Creates correlations between her results and Bob's
    \end{enumerate}
    
    \vspace{5mm}
    \begin{block}{Mathematical Description}
        Eve introduces local hidden variables $\mathbf{n}_a, \mathbf{n}_b$
        \begin{equation*}
            S = \int \rho(\mathbf{n}_a,\mathbf{n}_b) [\sqrt{2}\mathbf{n}_a\cdot\mathbf{n}_b] d\mathbf{n}_a d\mathbf{n}_b
        \end{equation*}
        This gives $-\sqrt{2} \leq S \leq \sqrt{2} < 2\sqrt{2}$
    \end{block}
\end{frame}

% ============ 第四部分：结果与分析 ============
\section{Results \& Analysis}

\begin{frame}{CHSH Test Results}
    \begin{table}[h]
        \centering
        \begin{tabular}{lccc}
            \toprule
            \textbf{Scenario} & \textbf{|S| Value} & \textbf{Quantum Limit?} & \textbf{Secure?} \\
            \midrule
            Ideal (no Eve) & 2.798 & Yes ($>2.0$) &  \yes \\
            With Eve Attack & 1.981 & No ($\leq 2.0$) &  \no \\
            \bottomrule
        \end{tabular}
        \caption{CHSH test results from simulation}
    \end{table}
    
    \vspace{3mm}
    \begin{columns}
        \begin{column}{0.5\textwidth}
            \textbf{Theoretical values:}
            \begin{itemize}
                \item Quantum: $|S| = 2\sqrt{2} \approx 2.828$
                \item Classical: $|S| \leq 2.000$
                \item Eve attack: $|S| \leq \sqrt{2} \approx 1.414$
            \end{itemize}
        \end{column}
        
        \begin{column}{0.5\textwidth}
            \begin{alertblock}{Security Threshold}
                $|S| > 2$: Secure quantum channel \\
                $|S| \leq 2$: Possible eavesdropping
            \end{alertblock}
        \end{column}
    \end{columns}
\end{frame}

\begin{frame}{Correlation Coefficients Comparison}
    \begin{figure}[h]
        \centering
        \includegraphics[width=0.8\textwidth]{correlation_plot.png}
        \caption{Correlation coefficients for different basis combinations}
    \end{figure}
    
    \begin{itemize}
        \item Perfect anticorrelation for matched bases: $E = -1$
        \item Quantum predictions closely matched in ideal case
        \item Eve's attack reduces correlation strength
    \end{itemize}
\end{frame}

\begin{frame}{Key Generation \& QBER Analysis}
    \begin{columns}
        \begin{column}{0.6\textwidth}
            \begin{table}[h]
                \centering
                \begin{tabular}{lcc}
                    \toprule
                    \textbf{Metric} & \textbf{Ideal} & \textbf{With Eve} \\
                    \midrule
                    Key bits available & 3,457 & 3,482 \\
                    Error bits & 12 & 867 \\
                    QBER & 0.35\% & 24.9\% \\
                    Key rate (bits/pair) & 0.35 & 0.35 \\
                    \bottomrule
                \end{tabular}
                \caption{Key generation performance}
            \end{table}
        \end{column}
        
        \begin{column}{0.4\textwidth}
            \begin{block}{QBER Threshold}
                Typical QBER limits:
                \begin{itemize}
                    \item BB84: $\sim 11\%$
                    \item E91: $\sim 14.6\%$
                    \item Our simulation: $24.9\%$ (Eve)
                \end{itemize}
            \end{block}
        \end{column}
    \end{columns}
    
    \vspace{3mm}
    \begin{alertblock}{Important Observation}
        Eve's attack increases QBER dramatically while keeping key rate similar
    \end{alertblock}
\end{frame}

% 合并后的单个frame（替代原来的2个frame）
\begin{frame}{Visualization of Results}
    % 1. 图片部分（调整宽度留出下方空间）
    \begin{figure}
        \centering
        % 宽度设为0.8\textwidth（比原来稍小，给列表留空间）
        \includegraphics[width=0.6\textwidth]{simulation_results.png}
        \caption{Complete simulation results visualization}
    \end{figure}
    
    \begin{enumerate}
        \item CHSH value comparison with theoretical limits
        \item Correlation coefficients for all basis combinations
        \item Available key bits count
        \item Quantum Bit Error Rate (QBER)
    \end{enumerate}
\end{frame}

% ============ 第五部分：讨论 ============
\section{Discussion}

\begin{frame}{Security Analysis}
    \textbf{Why E91 is secure:}
    \begin{itemize}
        \item No information in transit (delayed measurement principle)
        \item Bell test detects any eavesdropping
        \item Information "comes into being" only after public discussion
        \item No cloning theorem prevents copying of quantum states
    \end{itemize}
    
    \vspace{5mm}
    \textbf{Limitations in our simulation:}
    \begin{itemize}
        \item Simplified eavesdropper model
        \item No channel noise considered
        \item Perfect detection efficiency assumed
        \item Finite statistics effects
    \end{itemize}
\end{frame}

\begin{frame}{Comparison with BB84 Protocol}
    \begin{table}[h]
        \centering
        \begin{tabular}{lcc}
            \toprule
            \textbf{Feature} & \textbf{E91} & \textbf{BB84} \\
            \midrule
            Quantum resource & Entanglement & Single photons \\
            Security basis & Bell's theorem & No-cloning theorem \\
            Key extraction & From correlations & From bases matching \\
            Eavesdropping test & Statistical (CHSH) & Direct (QBER) \\
            Implementation complexity & Higher & Lower \\
            Experimental realization & More challenging & Well-established \\
            \bottomrule
        \end{tabular}
        \caption{E91 vs BB84 protocol comparison}
    \end{table}
    
    \vspace{3mm}
    \begin{block}{Advantage of E91}
        Provides device-independent security certification through Bell test
    \end{block}
\end{frame}

% ============ 第六部分：结论与展望 ============
\section{Conclusion \& Future Work}

\begin{frame}{Conclusion}
    \begin{itemize}
        \item \textbf{Successfully simulated} E91 quantum cryptography protocol
        \item \textbf{Verified} security based on Bell's theorem violation
        \item \textbf{Demonstrated} eavesdropping detection through CHSH test
        \item \textbf{Showed} practical key generation from quantum correlations
    \end{itemize}
    
    \vspace{5mm}
    \begin{block}{Key Findings}
        \begin{enumerate}
            \item Ideal case: $|S| = 2.798 > 2$ (secure)
            \item With Eve: $|S| = 1.981 \leq 2$ (insecure)
            \item QBER increases from 0.35\% to 24.9\% with eavesdropping
            \item Protocol successfully detects intercept-resend attacks
        \end{enumerate}
    \end{block}
\end{frame}

\begin{frame}{Future Work}
    \begin{columns}
        \begin{column}{0.5\textwidth}
            \textbf{Simulation improvements:}
            \begin{itemize}
                \item More realistic channel models
                \item Advanced eavesdropping strategies
                \item Finite-key analysis
                \item Device-independent aspects
            \end{itemize}
        \end{column}
        
        \begin{column}{0.5\textwidth}
            \textbf{Experimental directions:}
            \begin{itemize}
                \item Photon source optimization
                \item Detection efficiency improvements
                \item Real-time processing implementation
                \item Field deployment considerations
            \end{itemize}
        \end{column}
    \end{columns}
    
    \vspace{5mm}
    \begin{block}{Long-term Vision}
        \vspace{1mm}
        Towards practical device-independent quantum key distribution networks
    \end{block}
\end{frame}



\begin{frame}{}
    \centering
    \vspace{1cm}
    
    \Huge Thank You!
    
    \vspace{1cm}
    
    \large Questions \& Discussion
    
    \vspace{1.5cm}
    
    \normalsize
    Zhao Yumu \\
    School of Physical Science and Technology \\
    ShanghaiTech University \\
    \texttt{zhaoym2023@shanghaitech.edu.cn}
    
    \vspace{0.5cm}
    
    \footnotesize
    Code available at: \url{https://github.com/zhaomumu666/e91-simulation.git}
\end{frame}

% ============ 附录（可选） ============
\appendix

\begin{frame}{Appendix: Simulation Code Structure}
    \begin{columns}[T] % T: 分栏顶部对齐
        % 左栏：Key functions
        \column{0.5\textwidth}
        \textbf{Key functions:}
        \begin{itemize}
            \item \texttt{simulate\_ideal\_scenario()}
            \item \texttt{simulate\_eavesdropping\_scenario()}
            \item \texttt{calculate\_correlation()}
            \item \texttt{visualize\_results()}
            \item \texttt{main()} for execution
        \end{itemize}

        % 右栏：Main components
        \column{0.5\textwidth}
        \textbf{Main components:}
        \begin{itemize}
            \item \texttt{E91Simulator} class
            \item Quantum state preparation
            \item Measurement simulation
            \item Correlation calculation
            \item CHSH test implementation
            \item Visualization functions
        \end{itemize}
    \end{columns}

    % 中间加适度间距，分隔分栏与Technical Details
    \vspace{4mm}

    % 优化Technical Details的块样式，避免拥挤
    \begin{block}{Technical Details}
        \begin{itemize}
            \item Python 3.8+, Qiskit 0.45+
            \item 10,000 photon pairs simulated
            \item Reproducible random seeds
            \item Modular design for easy extension
        \end{itemize}
    \end{block}
    \end{frame}
\end{document}